\documentclass[aspectratio=169,11pt]{beamer}
\usetheme{default}
\setbeamertemplate{navigation symbols}{}
\setbeamertemplate{footline}{%
  \hfill\insertframenumber/\inserttotalframenumber\hspace*{4pt}\vspace{2pt}%
}
\usepackage{lmodern}
\usepackage[most]{tcolorbox}
\usepackage{listings}
\usepackage{tikz}
\usetikzlibrary{arrows.meta,positioning,shapes.geometric,fit,backgrounds}
\usepackage{booktabs}
\usepackage{hyperref}

\definecolor{lcgreen}{HTML}{2E7D32}
\definecolor{lcblue}{HTML}{1565C0}
\definecolor{lcorange}{HTML}{EF6C00}
\definecolor{lcred}{HTML}{C62828}
\definecolor{codebg}{HTML}{F5F5F5}
\definecolor{docgray}{gray}{0.55}

\newtcolorbox{conceptbox}[1][]{colback=yellow!20,colframe=yellow!60!black,title=#1,fonttitle=\bfseries,top=2pt,bottom=2pt,left=4pt,right=4pt}
\newtcolorbox{warningbox}[1][]{colback=red!15,colframe=red!60!black,title=#1,fonttitle=\bfseries,top=2pt,bottom=2pt,left=4pt,right=4pt}
\newtcolorbox{tipbox}[1][]{colback=green!10,colframe=green!50!black,title=#1,fonttitle=\bfseries,top=2pt,bottom=2pt,left=4pt,right=4pt}
\newtcolorbox{codebox}[1][]{colback=blue!10,colframe=blue!40!black,title=#1,fonttitle=\bfseries,top=2pt,bottom=2pt,left=4pt,right=4pt}

\lstset{
  basicstyle=\ttfamily\scriptsize,
  backgroundcolor=\color{codebg},
  breaklines=true,
  frame=single,
  rulecolor=\color{blue!40!black},
  keywordstyle=\color{lcblue}\bfseries,
  stringstyle=\color{lcgreen},
  commentstyle=\color{docgray}\itshape,
  language=Python,
  showstringspaces=false,
  tabsize=2,
  xleftmargin=2pt,
  xrightmargin=2pt,
  aboveskip=2pt,
  belowskip=2pt,
}

\newcommand{\docref}[1]{{\scriptsize\color{docgray}[Docs: #1]}}

\title{Putting It All Together}
\subtitle{Section 13: Building a Multi-Agent Research Assistant}
\author{LangGraph v1.0.x}
\date{March 2026}

\begin{document}

\frame{\titlepage}

\begin{frame}{Project Overview: Multi-Agent Research Assistant}
\setlength{\itemsep}{1pt}
\begin{itemize}
\item \textbf{Goal:} Build a collaborative agent system that researches topics and produces reports
\item \textbf{Agents:} Supervisor, Research, Writer, Reviewer
\item \textbf{Tools:} Web search, file I/O, LLM inference
\item \textbf{Features:} Streaming, persistence, human-in-the-loop, error handling
\item \textbf{Outcome:} Production-ready multi-agent system with all concepts integrated
\end{itemize}
\vspace{0.15cm}
\begin{conceptbox}[Integration Test]
This project validates state management, tools, subgraphs, persistence, streaming, and error handling.
\end{conceptbox}
\end{frame}

\begin{frame}{System Architecture}
\begin{center}
\begin{tikzpicture}[scale=0.85, every node/.style={font=\scriptsize}]
\node[draw,rounded rectangle,fill=blue!30] (input) at (0,4) {User Query};
\node[draw,rectangle,fill=purple!30] (supervisor) at (2,4) {Supervisor Agent};
\node[draw,rectangle,fill=green!30] (research) at (0,2) {Research Agent};
\node[draw,rectangle,fill=orange!30] (writer) at (2,2) {Writer Agent};
\node[draw,rectangle,fill=cyan!30] (reviewer) at (4,2) {Reviewer Agent};
\node[draw,rounded rectangle,fill=red!30] (output) at (2,0) {Final Report};
\node[draw,rectangle,fill=yellow!20] (tools) at (6,2.5) {Tools};
\draw[-{Latex}] (input) -- (supervisor);
\draw[-{Latex}] (supervisor) -- node[left] {delegate} (research);
\draw[-{Latex}] (supervisor) -- (writer);
\draw[-{Latex}] (supervisor) -- (reviewer);
\draw[-{Latex}] (research) -- (writer);
\draw[-{Latex}] (writer) -- (reviewer);
\draw[-{Latex}] (reviewer) -- (output);
\draw[-{Latex}] (research) -- node[right] {search, write} (tools);
\draw[-{Latex}] (writer) -- (tools);
\end{tikzpicture}
\end{center}
\docref{Multi-Agent Architecture}
\end{frame}

\begin{frame}{Project Structure}
\setlength{\itemsep}{1pt}
{\footnotesize
\begin{itemize}
\item \texttt{research\_assistant/}
\item \quad \texttt{main.py} - entry point, run graph
\item \quad \texttt{state.py} - shared state schema
\item \quad \texttt{agents/} - supervisor, research, writer, reviewer
\item \quad \texttt{tools/} - web search, file I/O
\item \quad \texttt{graphs/} - research subgraph, writer subgraph, etc.
\item \quad \texttt{config.py} - checkpointer, LLM settings
\end{itemize}
}
\vspace{0.1cm}
\begin{tipbox}[Organization]
Separate concerns: state, agents, tools, graphs. Clean separation enables testing and reuse.
\end{tipbox}
\end{frame}

\begin{frame}{Step 1: Define Shared State Schema}
\setlength{\itemsep}{1pt}
\begin{itemize}
\item Define: \texttt{class ResearchState(TypedDict)}
\item Fields: \texttt{query, research\_findings, draft\_report, feedback, final\_report, messages}
\item Use \texttt{Annotated[list, add]} for accumulating lists
\item Reducers automatically merge state updates across nodes
\end{itemize}
\vspace{0.1cm}
\begin{tipbox}[State Design]
Clear schema prevents confusion. Reducers ensure deterministic merging.
\end{tipbox}
\end{frame}

\begin{frame}{Step 2: Research Agent Subgraph with Tools}
\setlength{\itemsep}{1pt}
\begin{itemize}
\item Define tools: web search, file I/O, etc.
\item Bind to LLM: \texttt{model.bind\_tools([tool1, tool2])}
\item Invoke: \texttt{response = model\_with\_tools.ainvoke([...])}
\item Execute: if \texttt{response.tool\_calls}, run tools and return findings
\end{itemize}
\vspace{0.1cm}
\begin{tipbox}[Tool Integration]
Tools extend agent capabilities. LLM decides when to call them.
\end{tipbox}
\end{frame}

\begin{frame}{Step 3: Writer Agent Subgraph}
\setlength{\itemsep}{1pt}
\begin{itemize}
\item Input: research findings from previous step
\item Create prompt: "Based on findings... write a report for..."
\item Invoke LLM: \texttt{model.ainvoke([HumanMessage(prompt)])}
\item Return: \texttt{\{"draft\_report": response.content\}}
\end{itemize}
\vspace{0.1cm}
\begin{tipbox}[Agent Design]
Simple, focused agents work well. One task per node.
\end{tipbox}
\end{frame}

\begin{frame}{Step 4: Reviewer Agent Subgraph}
\setlength{\itemsep}{1pt}
\begin{itemize}
\item Input: draft report from writer
\item Prompt: "Review and suggest improvements: ..."
\item Invoke LLM for critical feedback
\item Return: \texttt{\{"feedback": response.content\}}
\end{itemize}
\vspace{0.1cm}
\begin{tipbox}[Quality Gate]
Reviewer acts as quality control. Feedback improves final output.
\end{tipbox}
\end{frame}

\begin{frame}{Step 5: Supervisor Graph (Routing)}
\setlength{\itemsep}{1pt}
\begin{itemize}
\item Create: \texttt{builder = StateGraph(ResearchState)}
\item Add nodes: research, writer, reviewer
\item Set entry: \texttt{builder.set\_entry\_point("research")}
\item Connect: research -> writer → reviewer → END
\end{itemize}
\vspace{0.1cm}
\begin{tipbox}[Linear Flow]
Simple linear pipeline. Can add conditional edges for retries or branches.
\end{tipbox}
\end{frame}

\begin{frame}{Step 6: Add Persistence (PostgreSQL)}
\setlength{\itemsep}{1pt}
\begin{itemize}
\item Import: \texttt{from langgraph.checkpoint.postgres import PostgresSaver}
\item Create: \texttt{checkpointer = PostgresSaver(conn\_string=...)}
\item Compile: \texttt{graph = builder.compile(checkpointer=checkpointer)}
\item Invoke with config: \texttt{config = \{"configurable": \{"thread\_id": ...\}\}}
\end{itemize}
\vspace{0.1cm}
\begin{tipbox}[Durability]
PostgreSQL enables crash recovery and multi-user state management.
\end{tipbox}
\end{frame}

\begin{frame}{Step 7: Human-in-the-Loop (Approval Gate)}
\setlength{\itemsep}{1pt}
\begin{itemize}
\item Add approval node: \texttt{builder.add\_node("approval", lambda s: s)}
\item Route to approval: \texttt{builder.add\_edge("reviewer", "approval")}
\item Set break: \texttt{builder.set\_break\_point("approval")}
\item Invoke: execution pauses at approval
\item Resume: \texttt{graph.invoke(None, config=config)}
\end{itemize}
\vspace{0.1cm}
\begin{tipbox}[Human Gates]
Enable workflow approval, manual corrections, or policy checks.
\end{tipbox}
\end{frame}

\begin{frame}{Step 8: Streaming Progress}
\setlength{\itemsep}{1pt}
\begin{itemize}
\item Use \texttt{graph.astream(input, config, stream\_mode="updates")}
\item Iterate: \texttt{async for event in graph.astream(...)}
\item Event contains: node name, data changes
\item Real-time feedback to UI, terminal, or logs
\end{itemize}
\vspace{0.1cm}
\begin{tipbox}[User Experience]
Streaming keeps users informed. No more "waiting for results".
\end{tipbox}
\end{frame}

\begin{frame}{Step 9: Error Handling \& Retries}
\setlength{\itemsep}{1pt}
\begin{itemize}
\item Import: \texttt{from langgraph.pregel import RetryPolicy}
\item Create: \texttt{RetryPolicy(max\_attempts=3, backoff\_factor=2.0, ...)}
\item Attach: \texttt{builder.add\_node(..., retry\_policy=rp)}
\item Auto-retries on specified exceptions with exponential backoff
\end{itemize}
\vspace{0.1cm}
\begin{tipbox}[Resilience]
Automatic retries handle transient failures without manual intervention.
\end{tipbox}
\end{frame}

\begin{frame}{Step 10: Testing the System}
\setlength{\itemsep}{1pt}
\begin{itemize}
\item Unit test: mock tools, verify node logic
\item Integration test: full graph with real config
\item Test happy path first
\item Add error injection: mock API failures, timeouts
\item Verify retry behavior, fallback activation
\end{itemize}
\vspace{0.1cm}
\begin{tipbox}[Test Coverage]
Test each node separately, then end-to-end. Errors caught early.
\end{tipbox}
\end{frame}

\begin{frame}{Step 11: Running the Application}
\setlength{\itemsep}{1pt}
\begin{itemize}
\item Define main coroutine with input state and config
\item Use \texttt{graph.astream(...)} for streaming
\item Iterate and process events
\item Call with \texttt{asyncio.run(main())}
\item Command: \texttt{python main.py}
\end{itemize}
\vspace{0.1cm}
\begin{tipbox}[Integration Complete]
All concepts working together: state, tools, streaming, persistence, error handling.
\end{tipbox}
\end{frame}

\begin{frame}{Concepts Used: Comprehensive Recap}
\small
\begin{tabular}{lll}
\toprule
\textbf{Section} & \textbf{Concept} & \textbf{Example} \\
\midrule
2 & StateGraph & State definition \\
3 & Nodes \& Edges & research -> writer → reviewer \\
4 & Tools & web search, file I/O \\
5 & Subgraphs & research subgraph \\
6 & Conditional Routing & route based on findings \\
7 & Human-in-the-Loop & approval gate \\
8 & Streaming & astream with updates \\
9 & Persistence & PostgresSaver \\
10 & Message History & message accumulation \\
12 & Error Handling & RetryPolicy + try/except \\
\bottomrule
\end{tabular}
\end{frame}

\begin{frame}{Key Lessons from Integration}
\setlength{\itemsep}{1pt}
\begin{itemize}
\item \textbf{Modularity:} Separate agents into nodes; compose into graph
\item \textbf{State clarity:} Define schema early; use Annotated for reducers
\item \textbf{Tool integration:} Agents use tools via model.bind\_tools()
\item \textbf{Robustness:} Combine retry policies, error handling, persistence
\item \textbf{Observability:} Stream events, log progress, monitor with LangSmith
\end{itemize}
\vspace{0.15cm}
\begin{tipbox}[Production Ready]
This architecture scales to 10+ agents, handles failures gracefully, and provides full observability.
\end{tipbox}
\end{frame}

\begin{frame}{Extending the System}
\setlength{\itemsep}{1pt}
\begin{itemize}
\item Add more specialist agents (fact-checker, summarizer, translator)
\item Use conditional edges to route based on content quality
\item Parallelize research across multiple queries
\item Add feedback loop: reviewer suggests research refinement
\item Deploy to LangGraph Platform for managed hosting
\end{itemize}
\vspace{0.15cm}
\begin{conceptbox}[Flexibility]
The modular structure enables easy addition of new agents without rewriting supervisor logic.
\end{conceptbox}
\end{frame}

\begin{frame}{Common Pitfalls}
\setlength{\itemsep}{1pt}
\begin{itemize}
\item \textbf{State explosion:} Too many fields -> hard to reason about reducer merges
\item \textbf{Tool chaos:} Agents with access to unrelated tools -> action confusion
\item \textbf{No async:} Synchronous agents block; use async for concurrency
\item \textbf{Missing error handlers:} Graph crashes on first tool failure
\item \textbf{No testing:} Assume integration works; test each agent separately first
\end{itemize}
\vspace{0.1cm}
\docref{Integration Best Practices}
\end{frame}

\begin{frame}{Summary: Building Multi-Agent Systems}
\setlength{\itemsep}{1pt}
\begin{enumerate}
\item Define state schema with Annotated reducers
\item Implement each agent as a node with clear inputs/outputs
\item Connect agents with edges (linear, conditional, parallel)
\item Add tools via model.bind\_tools()
\item Persist state with PostgresSaver
\item Enable human-in-the-loop via break points
\item Stream progress with astream()
\item Test each node independently before integration
\end{enumerate}
\vspace{0.15cm}
\docref{Multi-Agent Patterns}
\end{frame}

\begin{frame}{Self-Check Questions}
\setlength{\itemsep}{2pt}
\begin{enumerate}
\item How would you add a parallel research track (fetch from DB vs. web)?
\item What happens if the writer agent fails? How would you recover?
\item How would you add a feedback loop: reviewer -> research?
\item Why use PostgresSaver instead of in-memory state?
\item How would you stream partial results to a UI in real-time?
\item What metrics would you monitor in production?
\end{enumerate}
\end{frame}

\end{document}
